\section{Discussion}
\label{sec:discussion}

\subsection{What Decides Whether Anomaly Aiding Helps}

Across four observations, two estimators, and one consumer-grade IMU, the outcome is ordered by a single quantity: the quality of the anomaly measurement relative to the map signal. The two smartphone observations and the ADXL355, whose single-update scales along the map gradient are \boundscaleadxl{} to \boundscalephonemag~km, provide no benefit; the RM3100, whose scale of about \boundscalermthree~km is comparable to the error of the degraded solution, provides a consistent one. The inertial sensor was held fixed throughout, and the EKF and UKF behave alike, so neither the IMU nor the choice between the two Kalman filters is what limits the smartphone result. A better aiding sensor was sufficient to obtain a benefit; whether a better inertial sensor would add to it was not tested here.

The same heuristic gives an indicative requirement for these roads and maps. For a single update to resolve position along the gradient to the \bounddegradedkm~km of the degraded solution, the measurement noise must be below about \boundreqgrav~mGal for gravity and about \boundreqmag~nT for the magnetic field. Because the heuristic ignores the unobserved contour direction, the bias state, and the accumulation of repeated updates, this is an order-of-magnitude target rather than a threshold. A dedicated low-cost magnetometer approaches the magnetic requirement. No MEMS accelerometer approaches the gravity requirement; land-vehicle strapdown gravimetry has reached the milligal level only with navigation-grade inertial sensors, GNSS velocity, and filtering windows of minutes \citep{yu2015sins}. The map is also part of the heuristic. The WDMAM grid, at 3~arc-minutes and referenced to 5~km altitude, retains little structure at the sub-kilometer scale of the position error, and the free-air grid over land is effectively 9~km. Once the sensor is adequate, map resolution becomes the next limit, and higher-resolution regional surveys would lower it.

When the measurement noise greatly exceeds the map signal, the Kalman gain of the anomaly update falls toward zero and the filter effectively ignores the measurement, which is what both Kalman filters did with every observation except the RM3100. The particle filter has the analogous limit: a likelihood that is exactly independent of the state multiplies every weight by the same constant and leaves the normalized weights, and the systematic resampling of equal weights, unchanged. Its departures from that limit come from elsewhere, namely the finite sample of particles, the roughening jitter added after every resampling, and any mismatch between the modeled and actual measurement errors. The divergence of \Cref{sec:results-rbpf} occurs with and without aiding, so it does not show that weak anomaly updates harm the particle filter.

\subsection{The Particle Filter}

The results do not support the expectation, stated in the authors' conference paper, that a particle filter is the estimator best suited to anomaly aiding. Its theoretical advantage is the ability to represent a multimodal posterior, which arises when the position uncertainty spans several correlation lengths of the map. That is not the regime studied here. With a GNSS fix every minute, the median error of the Kalman-filter solution, about 0.44~km, is far smaller than the effective correlation lengths of the maps, about 9~km for gravity and about 30~km for the magnetic field, so the anomaly likelihood is effectively unimodal over the region of uncertainty. The particle filter's advantage should be sought in long outages, where the position uncertainty grows to the scale of the map features, and with a formulation that survives the degraded-GNSS baseline first. Two structural remedies for the deprivation of \Cref{sec:results-rbpf} are standard: applying GNSS position as a Kalman update of the full error state, which needs no particles because it is linear in position, and proposing particles near each fix, as in sensor resetting or a mixture proposal.

Two properties of the implementation would limit the benefit of aiding even without the divergence. The E\"{o}tv\"{o}s term makes the gravity observation depend on velocity, which the Rao-Blackwellized update does not model, so the likelihood can select particles by velocity rather than by position. And each filter treats successive anomaly errors as independent given the state. For the dedicated sensors this matches how the errors were generated, white noise plus a modeled Gauss--Markov bias over an exact map, so repeated observations of the same map feature do add information. For the smartphone observations it does not: the heading-dependent vehicle field that explains about half of the magnetic intensity variance (\Cref{sec:results-cause}) persists for as long as the heading does, so successive errors are correlated in time and the weights and covariances built on them are optimistic. The spatial correlation of the map alone does not imply correlated measurement errors.

\subsection{Limitations}
\label{sec:limitations}

Several limitations bound these conclusions. The reference for every error is the smartphone's own GNSS track; it limits the precision of small differences, but it applies equally to every configuration in a pair. The dedicated-sensor observations are idealized: the vehicle's magnetic field is absent and the maps are exact, so the magnetic improvement is an upper bound on what the same sensor would provide in a vehicle, where the platform field that dominated the smartphone magnetometer would have to be avoided or compensated. Their errors are a single realization, although each drive draws its own, so each paired test is over 27 independent draws. Anomaly errors are modeled as independent between the once-per-second updates; for the smartphone observations, whose vehicle-field error persists with heading, this is optimistic. The Kalman filters' measurement Jacobian omits the reference-field and motion terms of the observation; the direction of the resulting effect on navigation error is not established, and it is common to every configuration. The magnetometer heading observation shared by every run is itself corrupted by the vehicle's field, with errors far larger than its assumed 0.2~rad. Finally, the experiment covers one degradation scenario; a shorter fix interval would shrink the degraded error and, by the heuristic above, the benefit that any given sensor can provide, whereas a longer outage would enlarge both.
